% appendix.tex -- Appendix A: preregistration record (pins, audit record, families, deviations, symbol crosswalk).
%                 Appendix B: evidence that UniDisc's denoiser has no time input in the released configuration.
%                 Appendix D: proof of Proposition 2 (correlated-token gating; theory addendum r1, D0078; check log runs/P2-THEORY-CHECK).
% Sources: PREREG-GATE-v4-2026-10-28.md sections 1, 2, 4.1, 4.2, 4.4, 9; SAP-ADDENDUM-A4-2026-10-28.md (r1) A4.3.1, A4.5;
% REDTEAM-P2-PLAN.md section B; runs/P2-UNIDISC-KEYCHECK/README.md. Hashes are copied from the card, not recomputed here.

\section{Preregistration record}\label{app:prereg}

\subsection{Pins}\label{app:pins}
Table~\ref{tab:pins} lists the pinned artefacts.
Hashes of the card are SHA-256 of the committed bytes (LF convention) unless marked; the worktree is CRLF for some files, so a raw
hash can differ without a content change.
Twelve-hex prefixes stand for hashes that are given in full in the pin file of the repository.

\begin{table}[t]
\centering
\caption{Pinned artefacts of the main run. ``Pin file'' is \texttt{runs/P2-GATE-CARD-V4/pins.json}.}
\label{tab:pins}
\footnotesize
\begin{tabular}{p{0.24\linewidth}p{0.68\linewidth}}
\toprule
Artefact & Identifier \\
\midrule
Code, tests, guard (113 files) & commit \texttt{\seqsplit{9c4e1e25cb0d1a77b78410e1613098d8c6ba1085}} (final audited commit); moved by dated row A-1 to \texttt{\seqsplit{f2f7b2e15e4f2b90545837d9e64207051c24a1cd}}; per-file hashes in the pin file; manifest hash \texttt{\seqsplit{65513677ba4d98adfa1a861db4215200ad64b78aec1cceeaf0bcdca6b4664dec}} \\
Analysis plan with Addenda A1 to A3b.12 & \texttt{\seqsplit{e33f078c45b92b94cbf658192e2fd815f8c6eafbffe53ba39bc0dd728982afb7}} (LF) \\
Addendum A4 & draft commit \texttt{40828e96}, revision r1 commit \texttt{f1671c91} \\
Fold file & LF \texttt{\seqsplit{c19db5b3b9fabdeb89e8e686c58ae6c011f746529fbf9be84f1c58c99c446786}}; CRLF working copy \texttt{\seqsplit{450edb543779de7ee5e6a52422fd6507ef52f15ba3b956a8b05c2e561620f311}} \\
Hard-negative pool, manifest, leak check & \texttt{cfd0f94facd0} (pool, LF), \texttt{070fb29a2ebc} (manifest, LF), \texttt{91f4b7968b54} (leak check, LF) \\
Frozen fit files of the descriptive family E-pred & \texttt{c89700dbc3ba} (fit), \texttt{5c535a3bda75} (signal-to-noise calculator), \texttt{8d7528bbb37b} (test file); audit record \texttt{26145a2a6b91} \\
\bottomrule
\end{tabular}
\end{table}

\subsection{Audit and sign-off record}\label{app:audits}
The analysis plan was signed off by the statistician with changes.
The red team signed Addenda A2 and A3 with conditions, and signed Addendum A4 with eleven conditions that revision r1 answers.
The code was audited in several rounds by a separate auditor, who found defects that were fixed; the final verdict is a pass at the
final audited commit, and a delta audit covers the later change of two command-line flags.
The hard-negative id mapping and leak check were signed off by a data steward.
No signature says that the method works, and no reviewer read a GPU outcome.
% run: P2-A23-AUDIT
The pipeline audit that the positive control triggered (SAP A2.3) was carried out by a separate auditor on 29 October 2026, after the
outcome of PS was known; its verdict is no defect found, and its record is
\nolinkurl{research/paper2/prereg/A2.3-AUDIT-PS-2026-10-29.md} (data: \nolinkurl{runs/P2-A23-AUDIT/a23_ps_audit.json}).
Its descriptive account of the decision headroom is post hoc and is reported as such in Section~\ref{ssec:disc_findings}.

\subsection{Families and wording rules}\label{app:families}
Table~\ref{tab:families} lists the eleven test families and the three descriptive analyses.
\begin{table}[t]
\centering
\caption{Test families of the main run. $n$: planned number of items; $m_{\rm H}$: hypotheses in the family's Holm correction.
Each family has its own level 0.05; there is no control across families.}
\label{tab:families}
\footnotesize
\setlength{\tabcolsep}{3pt}
\begin{tabular}{p{0.12\linewidth}p{0.09\linewidth}p{0.36\linewidth}p{0.05\linewidth}p{0.28\linewidth}}
\toprule
Family & $n$ & Test & $m_{\rm H}$ & Not allowed \\
\midrule
P (headline) & 300, 500 or 300 with $d_R$ & exact McNemar, GATE-M against none: JPEG, noise, splice; under rule R3 Wilcoxon on $d_R$ & 3 & a claim from a rejection with $c>b$ in favour of the gate; ``predicted'', ``explained'', ``caused by detection'' \\
PS & as P & same test on patch and swap & 2 & ``the gate works where detection is easy''; use as a demonstration \\
S2 & 300 & exact McNemar I-GATE against I-RAND4 (kept draw), patch and splice; tested only if $n\ge120$ & 2 & ``improves'', ``restores'', gate detection as mechanism \\
B & 300 pairs & exact McNemar fusion against better single direction; untestable at ceiling 0.95 & 1 & a bound on foil noise; a statement about corruption \\
R & 150 & exact McNemar GATE-M against none per corruption of P$^+$; runs only if $\text{P}^+\neq\emptyset$ & $|\text{P}^+|$ & ``generalises''; pooling with UniDisc \\
S-M, S-Z, S-DN, S-D, S-D2, S-C & as P & mechanism, evidence ratio, denoising, detection, out-of-sample detection, clean control & 3, 3, 2, 2, 2, 1 & as in the protocol \\
\midrule
E-dec & P fold & paired bootstrap intervals only & -- & ``treatment'' or ``detection'' for its parts in abstract, highlights, conclusion \\
E-pred & P, PS cells & agreement of signs, Clopper--Pearson interval & -- & ``predicted'', ``explains'', ``supported'' \\
Suite & 100 + 100 & eight cells, intervals only & -- & any $p$-value in a claim; abstract or highlights \\
\bottomrule
\end{tabular}
\end{table}
No secondary family can rescue a null result in P.
A null in P is reported with its label and with no causal clause drawn from PS, S2, B, R, E-dec, E-pred, the suite or the
exploratory comparison of Section~\ref{ssec:tier3}, in the title, abstract, highlights or conclusion.

\subsection{Deviations and disclosures known at the time of writing}\label{app:deviations}
\begin{enumerate}
\item The sample size of S2 and B was set to 300 instead of 150 before any outcome, conditional on the first author's dated
confirmation; S2 was restricted to patch and splice after a 21-item dry run.
\item The first stage of the detection measurement stopped and was resumed at item 106 of 200.
\item Three checks of the second-model adapter are applied outside the module that holds the others; the adapter checks are executed
by the implementer's code and reviewed, not executed, by the auditor, so the adapter is not described as independently verified.
\item The sensitivity analyses of S2 and B read ``rejection'' as the favourable-direction rejection at the adjusted level; this can
over-flag sensitivity to $n$ and never under-flags it, and it enters no claim.
\item The firewall for the S2 smoke test is an honour rule.
\item The swap and patch parameters are stored in one frozen file that stands for two files named in the plan.
\item The chance bar for the sanity corruptions is called blind only on the lead's assertion that no detection-pilot output from item
26 onward was read, which nobody else can verify; the first author's own written co-attestation that he had not looked at any of those outputs is on file, dated before the analysis was run.
\item Where the card forbids any claim from a rejection with $c>b$ and the analysis plan prescribes the sentence ``reduced
verification accuracy'', the plan governs: no claim in favour of the gate is made from such a rejection, and no title carries it.
\item The exploratory comparison of Section~\ref{ssec:tier3} was pre-specified in Addendum A4, whose draft and revision were
committed before the headroom decision was written; the addendum's author states that no value of the headroom or primary stage was
read when it was written, which is an assertion and not an attestation.
Its code is committed only with a dated pin row.
\item Family B may run before S2 if a fixed time inequality holds; if a stop occurs while B runs, the order reverts, and B complete
with S2 partial is then possible.
\item Whether a fallback applies if P is null, and the wording rule that governs the framing of the paper in that case, was
\todo{AUTHOR: confirm the wording. Record (decisions D0076-D0077, CONFIRM-EVENING section 12): on 29 October 2026, about 00:10 IST, the
first author delegated the choice; the stricter reading of the outcome book was adopted before any value of the main test was read,
so the framing ``when does input gating help'' stays out of the title, abstract, highlights, conclusion and the opening of the
discussion, and the submission condition is unchanged with a status-only definition of ``informative''. Name who adopted it.}
% run: P2-GATE-M2 (chunks.jsonl, driver logs); operations record research/paper2/CONFIRM-EVENING-2026-10-28.md sections 10, 11, 13
\item The calibration stage (300 items) stopped twice on the slow-item rule of the guard, with 106 and then 114 rows complete: an item
took 3.0~s against a limit of twice a median of 1.5~s, and later 4.4~s against twice 1.8~s. Neither stop was a crash; the second coincided
with host-side load while the GPU waited. On the first author's dated answers the stage was relaunched, first unchanged and then with a
limit of three times the median item time and a floor of 2~s, and the same limit applies to every later stage except the second-model
stages (row A-2). Guard settings decide only whether a run continues; the rows already stored were kept unchanged.
% run: P2-GATE-M4 (driver logs)
\item The main test stage (500 items) stopped once at 334 of 500 on the default limit of 12 chunks of the stage runner (reason ``max
chunks reached''; not a guard stop and with no crash event) and was relaunched with a limit of 30 chunks, resuming at row 334 with the
same item list and the same flags. The preregistered launch command did not print this setting; it is an operational setting and changes
no output.
% run: P2-GATE-M5 (driver logs, telemetry)
\item During the coherence stage (100 items) the operating system restarted the laptop to install driver updates. The stage's process
ended after one of 100 items. The system log records an update-initiated restart, with no kernel-power, bug-check, hardware-error or GPU-driver
event, and the GPU driver version was unchanged, so this is not counted as a crash or a guard stop. The stage was continued afterwards on
the first author's instruction and completed with 100 of 100 items; automatic updates were paused afterwards.
\item The family order at the end of the schedule was FOIL before S2, which the plan allows if a fixed time inequality holds at the
end of the coherence stage (A4.3.2); it held. A separate auditor read the code of both stages (lines 1 to 149 of the FOIL runner among
others) and confirmed that no FOIL output can reach an S2 input. FOIL reads no output of the calibration, main-test, coherence, S2 or
E-pred stages; it reads, and may add, layout-only null files in a shared cache directory that S2 never reads. The scored path of FOIL uses
no random numbers: its only generators are local ones in the start-up checks and the bootstrap seed child used in its analysis. This is
a reading of code and not a run, and the rule that the analysis of FOIL waits until the S2 table exists is procedural and not enforced
by code.
\item No second-model stage was run. The set P$^+$ of corruptions with a primary rejection and $b>c$ is empty, and by the plan
(AD.7) family R is run only for P$^+$.
% run: P2-S2-SELFCHECK-DIAG (figures from research/paper2/method/S2-SELFCHECK-DIAG-2026-10-29.md and card v4 row A-5)
\item The smoke test of S2 stopped twice at the tokenisation self-check, before any item, because 0.9854 and then 0.9844 of the image
tokens matched the reference where the rule requires 0.99 (the main-test stage gave 1.0). The cause was process-to-process
nondeterminism of the autotuned TF32 convolutions in the image tokeniser: the same image gave 0.998 in one fresh process and 0.985 in
another, with no S2 code involved. The self-check was changed to run with cuDNN autotuning off and deterministic algorithms on, with the
previous settings restored afterwards; the 0.99 rule, the scoring and the sampler are unchanged (dated row A-5, code commit
\texttt{8d1bee3a}, passed a delta audit). The check then passed. S2 reads the stored codes of the main-test stage and never tokenises an
image afresh, so no S2 output depended on the check. The directories of the two failed smoke attempts are kept.
\end{enumerate}

\subsection{Dated amendments}\label{app:amend}
Row A-1 moved the code pin to the commit named in Table~\ref{tab:pins} for every stage launched afterwards; it changed two
command-line flags of the stage runner and the guard settings recorded per chunk, and no analysis path, scoring definition, frozen
file or guard.
Row A-2 appended a slow-item limit of three times the median item time, with a floor of two seconds, to every non-second-model
stage, on the first author's dated answer; guard settings decide only whether a run continues and never change an item's output.
Row A-5 moved the code pin to commit \texttt{8d1bee3a} for the S2 stage and later stages; it changed only the S2 runner, so that its
unchanged self-check runs with deterministic convolution algorithms, and added one test (deviation list above).

\subsection{Symbol crosswalk}\label{app:crosswalk}
Table~\ref{tab:crosswalk} maps the symbols of the preregistration documents and code to the paper's.
\begin{table}[t]
\centering
\caption{Crosswalk from preregistration and code names to the symbols of this paper.}
\label{tab:crosswalk}
\footnotesize
\begin{tabular}{p{0.38\linewidth}p{0.5\linewidth}}
\toprule
Preregistration or code & This paper \\
\midrule
\texttt{alpha}, \texttt{beta} in the signal-to-noise calculator & $\phi_0$, $\phi_1$ \\
\texttt{alpha\_obs}, \texttt{beta\_obs} & $\hat\phi_0$, $\hat\phi_1$ \\
CRC level $\alpha$ & $\alpha$ \\
partition group $S_k$ & $\Gamma_k$ \\
re-infill samples (\texttt{n\_draw}, $L$) & $n_{\rm draw}$ \\
I-GATE, I-RAND4 & GATE-I, random-gate best-of-four re-infill \\
Proposition 4 of the specification (plug-in signal-to-noise change) & Eq.~\eqref{eq:snr} \\
Proposition 1 of the paper (conformal risk control applied) & Proposition~\ref{prop:crc} \\
\bottomrule
\end{tabular}
\end{table}

\section{Time input of the released configuration}\label{app:timefree}
Section~\ref{ssec:data} states that UniDisc's denoiser has no time input in the released configuration.
Table~\ref{tab:timefree} gives the evidence, read in the authors' code at commit
\texttt{\seqsplit{3ee3e47a9e64bac2fdca59957726f14723371ced}} of the UniDisc repository (clean tree) and in the key names of the
checkpoint, snapshot \texttt{\seqsplit{2b344ca856bd108b52a2d022a3463dace131249f}}.
The statement is about this configuration and checkpoint and not about UniDisc in general.
\begin{table}[t]
\centering
\caption{Evidence for the absence of a time input. Paths are relative to the UniDisc repository unless marked.}
\label{tab:timefree}
\footnotesize
\begin{tabular}{p{0.4\linewidth}p{0.5\linewidth}}
\toprule
Claim & Evidence (file:line) \\
\midrule
Time conditioning is off; the diffusion type is absorbing state, the parameterisation \texttt{subs}; no experiment overrides it & \texttt{configs/config.yaml:12,14,15} \\
The project composes the official demonstration experiments and loads the interleaved checkpoint & project files \texttt{runs/env-2026-10-14/\allowbreak scripts/\allowbreak unidisc\_\allowbreak setup.py:17,30-35} and \texttt{exp001/\allowbreak unidisc\_\allowbreak net.py:46-47} \\
The official inference experiment sets a null noise level; the authors' samplers pass \texttt{None} & \texttt{configs/\allowbreak experiments/\allowbreak large\_\allowbreak scale\_\allowbreak high\_\allowbreak res\_\allowbreak interleaved\_\allowbreak inference.\allowbreak yaml:30-31}; \texttt{model\_eval.py:\allowbreak 2057,\allowbreak 2086,\allowbreak 3019,\allowbreak 3060,\allowbreak 3186} \\
The forward pass accepts a null noise level only when this is allowed & \texttt{model.py:660-662,691,729} \\
No timestep embedder and no time modulation in the final layer without time conditioning & \texttt{models/\allowbreak dit.py:\allowbreak 1105,\allowbreak 1186-1188,\allowbreak 1377-1379} and \texttt{models/\allowbreak dit.py:\allowbreak 1064-1090} \\
A stored scaling flag is never read & \texttt{configs/\allowbreak model/\allowbreak extra\_\allowbreak large.yaml:8}; \texttt{models/\allowbreak dit.py:1290} \\
The checkpoint has no time-embedding or time-modulation keys; the string ``sigma'' does not occur in it & key names read from the checkpoint pickle, no tensor loaded \\
Raw logits are returned before the step that copies observed positions & \texttt{model.py:784-788} \\
\bottomrule
\end{tabular}
\end{table}

The loader is non-strict by default.
% run: P2-UNIDISC-KEYCHECK
A key check on a model built without weights found 342 of 342 checkpoint tensors matched to model parameters with equal shape and
dtype, no missing key, no unexpected key, and no shape mismatch, and the framework's own loader reports empty lists of missing and
unexpected keys.
Deterministic rotary-embedding buffers are computed at construction and are not part of any state dictionary.
The check covers the denoiser backbone, not the image tokeniser, which is loaded by another path, and it does not show that the
checkpoint values are the intended final weights.
During training, a whole modality was masked with some probability; which value applied to the released checkpoint is not recoverable
from the code.

\section{Diagnostics, cost and hardware}\label{app:diag}
\subsection{Calibration diagnostics and pipeline checks}\label{app:diagchecks}
On the clean test items the realised false-alarm proportion $\hat F=n^{-1}\sum_i L_i(\hat\tau)$ is reported for
$\alpha\in\{0.01,0.05,0.10\}$, with $\alpha=0.05$ primary.
A violation is flagged when $Z_{\rm cal}=(\hat F-\alpha)/\big(\hat\sigma_L\sqrt{1/n_{\rm cal}+1/n_{\rm test}}\big)$ exceeds 2.13 at any
level, where $\hat\sigma_L$ is the pooled item-level standard deviation of $L_i$; a pass means ``no evidence of violation'' and is
never worded as validation.
The precision among flagged tokens, the false-alarm proportion on the clean tokens of corrupted items and the spatial map of clean
false alarms are reported with equal prominence; for $\alpha=0.05$ on an $8\times8$ patch image, about 48 false flags are expected
against at most 64 true ones, so precision is at most 0.57 even for a perfect detector.
For the sanity corruptions the lower bound of the 95\,\% interval of the per-item AUROC of $q$ and of plain-$B$ must exceed 0.5; a
result at chance is a pipeline defect, not a scientific result, and triggers an audit.
The clean no-gate accuracy is reported with a Wilson interval against the pilot value of Section~\ref{ssec:pilot} as a descriptive
check.
A coherence test on 100 clean calibration items measures whether the denoiser's joint likelihood of two image tokens is independent
of the order in which they are revealed; it is read as ``coherent within numerical noise'' if the median absolute gap lies within the
95th percentile of the gap between repeats of the same computation, and as ``incoherent at the evidence scale'' if it exceeds the
clean mean per-token paired evidence, in which case Eq.~\eqref{eq:marginalise} is reported as not applicable.

\subsection{Cost in network calls}\label{app:cost}
Cost is counted in network passes (one row of a batch), because a call at this scale is the dominant term.
Table~\ref{tab:cost} lists the passes per quantity.
Gating the image adds $K$ passes per image version and nothing else to the image-direction evidence, because the evidence
passes depend on the caption only.
With only the image gated, verification of $N$ candidates costs $K$ gate passes, $2N$ evidence passes (the image-direction
half cached per caption), and $S_rn_{\rm draw}$ more for GATE-I, against $2N$ for the ungated two-way baseline.

\begin{table}[t]
\centering
\caption{Network passes per quantity. $K$: complementary groups; $N$: candidate captions; $S_r$, $n_{\rm draw}$: re-infill steps and draws.}
\label{tab:cost}
\small
\begin{tabular}{p{0.36\linewidth}p{0.12\linewidth}p{0.38\linewidth}}
\toprule
Quantity & Passes & Depends on, and cache \\
\midrule
$B$ over all image tokens & $K$ & image only (caption masked); per image \\
$C$ for image tokens (image fully masked) & 1 & caption only; per caption, shared across images and arms \\
$D$ for image or caption tokens & 1 & layout only (caption length); per length \\
$A$ over all image tokens & $K$ & image and caption \\
Image-direction evidence, all arms & 2 & the two captions; arms differ only in gathered entries \\
Text-direction evidence per candidate and arm & 1 & image, gated set \\
Re-infill (GATE-I) & $S_rn_{\rm draw}$ & image, gated set; shared by all candidates \\
\bottomrule
\end{tabular}
\end{table}

\subsection{Hardware and runtime}\label{app:hardware}
% run: P2-PROBE-01
All runs use one laptop GPU (NVIDIA RTX 5070 Ti Laptop, 12~GB) with 31~GB of system memory and no other hardware.
In the pilot one item cost 58 network passes, the median item time was 7.0~s and the peak GPU memory 8.2~GB, so passes are batched
four rows at a time.
Runs go in short checkpointed chunks under a guard that stops on memory, temperature and power-source events and is resumable; an
abort ends the stage and is not retried until reviewed.
% run: P2-PROBE-02
An earlier detection pilot stopped after 20 complete items when the GPU session failed.
% run: P2-S2-SELFCHECK-DIAG
For reproducibility: image tokenisation in a fresh process runs with autotuned TF32 convolutions, and two processes can differ by up to
about 1.5\,\% of the 1\,024 tokens of an image (self-check agreement 0.998 to 1.0 in earlier stages, 0.984 and 0.985 in the S2 smoke
test); the stored codes of the main-test stage are the reference, and S2 reads them and does not tokenise images afresh.

\section{Proof of Proposition~\ref{prop:snrcorr}}\label{app:snrcorr}
Write $m_0=(1-\rho)m$ and $m_1=\rho m$ for the numbers of clean and corrupted tokens, $k_0=1-\phi_0$ and $k_1=1-\phi_1$ for the keep
probabilities, and $z_j\in\{0,1\}$ for the label of token $j$, so that $\E\Delta_j$ is $\mu$ or $\mu_c$ and $\Var\Delta_j$ is $\sigma^2$ or
$\sigma_c^2$.
The exchangeable block covariance of assumption (a) is positive semi-definite if and only if $\gamma_0\ge-1/(m_0-1)$,
$\gamma_1\ge-1/(m_1-1)$, $\gamma_0,\gamma_1\le1$ and $\gamma_{01}^2\,m_0m_1\le\{1+(m_0-1)\gamma_0\}\{1+(m_1-1)\gamma_1\}$.

\paragraph{Variance.}
By assumption (b), $I_j$ is independent of $\Delta_j$, so
$\Var(I_j\Delta_j)=k(\sigma_z^2+\mu_z^2)-k^2\mu_z^2=k\sigma_z^2+k(1-k)\mu_z^2$ with $k=k_{z_j}$, which are the terms of Eq.~\eqref{eq:var}.
For $j\neq k$, again by (b),
$\mathrm{Cov}(I_j\Delta_j,I_k\Delta_k)=\E I_j\,\E I_k\,\E[\Delta_j\Delta_k]-\E I_j\,\E I_k\,\E\Delta_j\,\E\Delta_k=k_{z_j}k_{z_k}\,\mathrm{Cov}(\Delta_j,\Delta_k)$.
There are $m_0(m_0-1)$ ordered clean pairs with covariance $\gamma_0\sigma^2$, $m_1(m_1-1)$ ordered corrupted pairs with covariance
$\gamma_1\sigma_c^2$ and $2m_0m_1$ ordered mixed pairs with covariance $\gamma_{01}\sigma\sigma_c$, which gives Eq.~\eqref{eq:varcorr}.

\paragraph{Condition.}
Let $V_g$ be Eq.~\eqref{eq:varcorr} at $(\phi_0,\phi_1)$ and $V_0$ its value at $\phi_0=\phi_1=0$.
With $\mu_c=0$, $\E S=m_0k_0\mu$ with the gate and $m_0\mu$ without it, so, both being positive, the gate raises the ratio if and only if
$k_0^2V_0>V_g$.
Collect $k_0^2V_0-V_g$ by type of term.
The clean terms are
$k_0^2m_0\sigma^2\{1+(m_0-1)\gamma_0\}-m_0k_0(\sigma^2+\phi_0\mu^2)-m_0(m_0-1)k_0^2\gamma_0\sigma^2=-m_0k_0\phi_0(\sigma^2+\mu^2)$,
in which the $\gamma_0$ terms cancel.
The corrupted terms are $m_1\sigma_c^2\{(k_0^2-k_1)+(m_1-1)\gamma_1(k_0^2-k_1^2)\}$, with $k_0^2-k_1=\phi_1-2\phi_0+\phi_0^2$ and
$k_0^2-k_1^2=(\phi_1-\phi_0)(2-\phi_0-\phi_1)$.
The mixed terms are $2m_0m_1\gamma_{01}\sigma\sigma_c\,k_0(k_0-k_1)$, with $k_0-k_1=\phi_1-\phi_0$.
Dividing the sum by $m$ gives Eq.~\eqref{eq:snrcorr}.

\paragraph{Consequences.}
(i) is read off Eq.~\eqref{eq:snrcorr}.
(ii) At $\phi_1=\phi_0\in(0,1)$ every correlation term vanishes and the left side is $\rho\sigma_c^2(\phi_0^2-\phi_0)<0$, below the positive
right side.
For $\phi_1\le\phi_0$, let $G=(\rho m-1)\gamma_1$, so that $G\ge-1$ by positive semi-definiteness.
The bracket $(\phi_1-2\phi_0+\phi_0^2)+G(\phi_1-\phi_0)(2-\phi_0-\phi_1)$ is linear in $G$; it equals
$\phi_1-2\phi_0+\phi_0^2\le-\phi_0(1-\phi_0)$ at $G=0$, equals $-\phi_1(1-\phi_1)\le0$ at $G=-1$, and does not increase for $G\ge0$; hence it is
at most $0$.
If $\gamma_{01}\ge0$ the mixed term is at most $0$ as well, while the right side is positive for $\phi_0\in(0,1)$.
(iii) The derivatives of the left side with respect to $\gamma_1$ and $\gamma_{01}$ are
$\rho\sigma_c^2(\rho m-1)(\phi_1-\phi_0)(2-\phi_0-\phi_1)$ and $2\rho(1-\rho)m\sigma\sigma_c(1-\phi_0)(\phi_1-\phi_0)$, both positive when
$\phi_1>\phi_0$ and $\rho m>1$, and the right side does not depend on $\gamma_1$ or $\gamma_{01}$.
For fixed $\phi_1>\phi_0$ the $\gamma_1$ term grows without bound with $(\rho m-1)\gamma_1$, which gives the limit.
With $\gamma_{01}=0$, $x=\phi_1-\phi_0$ and $\chi=\rho\sigma_c^2/\{(1-\rho)(\sigma^2+\mu^2)\}$, Eq.~\eqref{eq:snrcorr} reads
$Gx^2-x\{1+2G(1-\phi_0)\}+\phi_0(1-\phi_0)(1+1/\chi)<0$, so for $G>0$ the gate raises the ratio if and only if $\phi_1$ exceeds
$\phi_0+\big[1+2G(1-\phi_0)-\sqrt{\{1+2G(1-\phi_0)\}^2-4G\phi_0(1-\phi_0)(1+1/\chi)}\big]/(2G)$ (and $\phi_1\le1$); as $G\to0$ this tends to the
boundary $2\phi_0-\phi_0^2+\phi_0(1-\phi_0)/\chi$ of Eq.~\eqref{eq:snr}.
(iv) For $\gamma_{01}<0$ and $\phi_1<\phi_0$ the mixed term is positive and grows with $m$, so it can exceed the remaining terms.
% run: P2-THEORY-CHECK (check D4)
An instance with a positive definite covariance is $m_0=900$, $m_1=100$, $\gamma_0=\gamma_1=0.05$, $\gamma_{01}=-0.045$, $\mu=0.1$, $\mu_c=0$,
$\sigma^2=\sigma_c^2=1$, for which the smallest eigenvalue of the correlation matrix is 0.9500 and the ratio rises from 0.4892 without a
gate to 0.5080 with the gate $(\phi_0,\phi_1)=(0.3,0)$; the instance is an illustration of the algebra and not an estimate from the model.
(v) $m$ enters Eq.~\eqref{eq:snrcorr} through $\rho m-1$ and $m$ whenever $\gamma_1\neq0$ or $\gamma_{01}\neq0$. \qed
